\chapter*{ЗАКЛЮЧЕНИЕ}
\addcontentsline{toc}{chapter}{ЗАКЛЮЧЕНИЕ}

В ходе данной работы были выполнены следующие задачи:
\begin{itemize}
    \item проведен анализ предметной области предметной области рекомендательных систем;
    \item формализована задача составления рекомендаций;
    \item проведен обзор существующих методов или групп методов решения;
    \item сформулированы критерии сравнения алгоритмов рекомендательных систем;
    \item выполнено сравнение перечисленные методов по сформулированным критериям.
\end{itemize}

Была достигнута поставленная цель данной работы.
В ходе исследования был проведен анализ основных типов рекомендательных систем,
включающий коллаборативную фильтрацию, контентно-ориентированные методы, методы,
основанные на знаниях, и гибридные подходы.

Результаты сравнения показали, что каждый тип систем имеет свои преимущества и недостатки.
Коллаборативная фильтрация наиболее эффективна при наличии большого объема данных о
взаимодействиях пользователей, но испытывает трудности с проблемой холодного старта.
Контентно-ориентированные системы хорошо работают в условиях ограниченной информации о
пользователях, но их эффективность снижается при отсутствии качественных атрибутов объектов.
Методы, основанные на знаниях, являются интерпретируемыми и подходят для доменов,
где требуется высокая точность рекомендаций, но требуют значительных усилий для построения базы
знаний. Гибридные системы сочетают сильные стороны различных подходов, обеспечивая универсальность,
разнообразие и масштабируемость рекомендаций, однако их внедрение требует значительных вычислительных ресурсов.
